% Options for packages loaded elsewhere
\PassOptionsToPackage{unicode}{hyperref}
\PassOptionsToPackage{hyphens}{url}
\documentclass[
]{article}
\usepackage{xcolor}
\usepackage[margin=1in]{geometry}
\usepackage{amsmath,amssymb}
\setcounter{secnumdepth}{-\maxdimen} % remove section numbering
\usepackage{iftex}
\ifPDFTeX
  \usepackage[T1]{fontenc}
  \usepackage[utf8]{inputenc}
  \usepackage{textcomp} % provide euro and other symbols
\else % if luatex or xetex
  \usepackage{unicode-math} % this also loads fontspec
  \defaultfontfeatures{Scale=MatchLowercase}
  \defaultfontfeatures[\rmfamily]{Ligatures=TeX,Scale=1}
\fi
\usepackage{lmodern}
\ifPDFTeX\else
  % xetex/luatex font selection
\fi
% Use upquote if available, for straight quotes in verbatim environments
\IfFileExists{upquote.sty}{\usepackage{upquote}}{}
\IfFileExists{microtype.sty}{% use microtype if available
  \usepackage[]{microtype}
  \UseMicrotypeSet[protrusion]{basicmath} % disable protrusion for tt fonts
}{}
\makeatletter
\@ifundefined{KOMAClassName}{% if non-KOMA class
  \IfFileExists{parskip.sty}{%
    \usepackage{parskip}
  }{% else
    \setlength{\parindent}{0pt}
    \setlength{\parskip}{6pt plus 2pt minus 1pt}}
}{% if KOMA class
  \KOMAoptions{parskip=half}}
\makeatother
\usepackage{longtable,booktabs,array}
\newcounter{none} % for unnumbered tables
\usepackage{calc} % for calculating minipage widths
% Correct order of tables after \paragraph or \subparagraph
\usepackage{etoolbox}
\makeatletter
\patchcmd\longtable{\par}{\if@noskipsec\mbox{}\fi\par}{}{}
\makeatother
% Allow footnotes in longtable head/foot
\IfFileExists{footnotehyper.sty}{\usepackage{footnotehyper}}{\usepackage{footnote}}
\makesavenoteenv{longtable}
\usepackage{graphicx}
\makeatletter
\newsavebox\pandoc@box
\newcommand*\pandocbounded[1]{% scales image to fit in text height/width
  \sbox\pandoc@box{#1}%
  \Gscale@div\@tempa{\textheight}{\dimexpr\ht\pandoc@box+\dp\pandoc@box\relax}%
  \Gscale@div\@tempb{\linewidth}{\wd\pandoc@box}%
  \ifdim\@tempb\p@<\@tempa\p@\let\@tempa\@tempb\fi% select the smaller of both
  \ifdim\@tempa\p@<\p@\scalebox{\@tempa}{\usebox\pandoc@box}%
  \else\usebox{\pandoc@box}%
  \fi%
}
% Set default figure placement to htbp
\def\fps@figure{htbp}
\makeatother
% definitions for citeproc citations
\NewDocumentCommand\citeproctext{}{}
\NewDocumentCommand\citeproc{mm}{%
  \begingroup\def\citeproctext{#2}\cite{#1}\endgroup}
\makeatletter
 % allow citations to break across lines
 \let\@cite@ofmt\@firstofone
 % avoid brackets around text for \cite:
 \def\@biblabel#1{}
 \def\@cite#1#2{{#1\if@tempswa , #2\fi}}
\makeatother
\newlength{\cslhangindent}
\setlength{\cslhangindent}{1.5em}
\newlength{\csllabelwidth}
\setlength{\csllabelwidth}{3em}
\newenvironment{CSLReferences}[2] % #1 hanging-indent, #2 entry-spacing
 {\begin{list}{}{%
  \setlength{\itemindent}{0pt}
  \setlength{\leftmargin}{0pt}
  \setlength{\parsep}{0pt}
  % turn on hanging indent if param 1 is 1
  \ifodd #1
   \setlength{\leftmargin}{\cslhangindent}
   \setlength{\itemindent}{-1\cslhangindent}
  \fi
  % set entry spacing
  \setlength{\itemsep}{#2\baselineskip}}}
 {\end{list}}
\usepackage{calc}
\newcommand{\CSLBlock}[1]{\hfill\break\parbox[t]{\linewidth}{\strut\ignorespaces#1\strut}}
\newcommand{\CSLLeftMargin}[1]{\parbox[t]{\csllabelwidth}{\strut#1\strut}}
\newcommand{\CSLRightInline}[1]{\parbox[t]{\linewidth - \csllabelwidth}{\strut#1\strut}}
\newcommand{\CSLIndent}[1]{\hspace{\cslhangindent}#1}
\setlength{\emergencystretch}{3em} % prevent overfull lines
\providecommand{\tightlist}{%
  \setlength{\itemsep}{0pt}\setlength{\parskip}{0pt}}
\usepackage{subfig}
\usepackage{setspace}\onehalfspacing
\usepackage{amsmath}
\usepackage{amsfonts}
\usepackage{amsthm}
\usepackage{mitpress}
\newtheorem{prop}{Proposition}
\newtheorem{claim}{Claim}
\usepackage{bookmark}
\IfFileExists{xurl.sty}{\usepackage{xurl}}{} % add URL line breaks if available
\urlstyle{same}
\hypersetup{
  pdftitle={Model section},
  hidelinks,
  pdfcreator={LaTeX via pandoc}}

\title{Model section}
\author{}
\date{\vspace{-2.5em}}

\usepackage{amsthm}
\newtheorem{theorem}{Theorem}
\newtheorem{lemma}{Lemma}
\newtheorem{corollary}{Corollary}
\newtheorem{proposition}{Proposition}
\newtheorem{conjecture}{Conjecture}
\theoremstyle{definition}
\newtheorem{definition}{Definition}
\theoremstyle{definition}
\newtheorem{example}{Example}
\theoremstyle{definition}
\newtheorem{exercise}{Exercise}
\theoremstyle{definition}
\newtheorem{hypothesis}{Hypothesis}
\theoremstyle{remark}
\newtheorem*{remark}{Remark}
\newtheorem*{solution}{Solution}
\begin{document}
\maketitle

People in the marriage market have two characteristics:
\(x=\left( x_{1},x_{2}\right)\), drawn from a normal distribution
\[
\mathcal{N}
\left( 
\begin{array}{c}
0 \\ 
0%
\end{array}%
,%
\begin{array}{cc}
s^{2} & \sigma \\ 
\sigma & S^{2}%
\end{array}%
\right).
\]
We interpret \(x_1\) as a genetic measure, for example of genetic variants predicting
height, physical attractiveness, health or intelligence. \(x_2\) is a measure of
socio-economic status, such as income or wealth, or social status more generally
(we sometimes use ``wealth'' as a shorthand).

People's attractiveness is given by
\[
i\left( x\right) =ax_{1}+\left( 1-a\right) x_{2}
\]
where \(a \in [0, 1]\) is a parameter reflecting the relative importance of
genetics to wealth in the marriage market.\footnote{Note that since the variance of the shocks to \(x_1\) and \(x_2\)
  (see below) has been normalized to 1, \(a\) also reflects this variance. That
  is, a large variance of SES shocks (compared to genetic shocks) translates into
  \(a\) being large.} If \(a = 0\),
marriage markets are highly inegalitarian, such that only SES matters. If \(a =
1\), marriage markets are economically egalitarian and only genetics matter. We
expect realistic societies to fall between these extremes, with \(0 < a < 1\).
Then, both genes and SES matter to attractiveness, and as a result,
social-genetic assortative mating (SGAM) takes place. This model does not assume
that people match \emph{directly} on genetics, which most observers agree would be
unlikely. Instead, genetics and SES both contribute to an attractiveness phenotype,
which is matched on. Attractiveness is a deterministic function of genes and SES,
but both genes and SES have random shocks in each generation (see below). Adding
a random component to attractiveness itself would likely retain the basic insights
of our model whilst reducing the effects of SGAM on genes-SES correlation.

Attractiveness \(i\) is
distributed \(N(0,\sigma_{I}^{2})\), where
\[
\sigma _{I}^{2}=a^{2}s^{2}+\left( 1-a\right) ^{2}S^{2}+2a\left( 1-a\right)\sigma.
\]

People form matches with transferable utility, where the surplus for a match
between \(x\) and \(y\) is \(S(i(x), i(y))\) such that \(\partial^{2}S/\partial
i\partial j > 0\), i.e.~\(S\) is supermodular. As a result there is positive
assortative mating on attractiveness: \(x\) matches with \(y\) only if they are at
the same quantile of attractiveness, i.e.~if \(i(x_{1},x_{2}) = i(y_{1},y_{2})\).
Within attractiveness quantiles, matching is random. This is the
SGAM mechanism.

Our first proposition shows that if SGAM is taking place, i.e.~if \(0 < a < 1\), then
there is a positive correlation between one partner's wealth and the other
partner's genetics.

\begin{proposition}\label{prop-couples-SGAM}
Under SGAM, the distribution of couples' characteristics is normal, with mean 0
and covariance matrix
\begin{equation}\label{cov-couples-SGAM}
\mathbb{C}\left( 
\begin{array}{c}
x_{1} \\ 
x_{2} \\ 
y_{1} \\ 
y_{2}%
\end{array}%
\right) =\allowbreak \left( 
\begin{array}{cccc}
s^{2} & \sigma  & A^{2} & AC \\ 
\sigma  & S^{2} & AC & C^{2} \\ 
A^{2} & AC & s^{2} & \sigma  \\ 
AC & C^{2} & \sigma  & S^{2}%
\end{array}%
\right) \allowbreak 
\end{equation}
where:\footnote{$A$ and $C$ are not related to the ACE model in genetics.}
\begin{align*}
A &= \frac{as^{2}+\left( 1-a\right) \sigma }{\sqrt{a^{2}s^{2}+\left(
1-a\right) ^{2}S^{2}+2a\left( 1-a\right) \sigma }} &= \frac{as^{2}+\left( 1-a\right) \sigma }{\sigma_I}; \\
C &= \frac{a\sigma +\left( 1-a\right) S^{2}}{\sqrt{a^{2}s^{2}+\left(
1-a\right) ^{2}S^{2}+2a\left( 1-a\right) \sigma }} &= \frac{a\sigma +\left( 1-a\right) S^{2}}{\sigma_I}.
\end{align*}

In particular, the covariance between $x_2$ and $y_1$, $AC$, is positive if
either $x_1$ and $x_2$ are already correlated ($\sigma > 0$), or if they are
uncorrelated ($\sigma = 0$) and the attractiveness parameter $a$ is strictly
between 0 and 1.

\begin{proof}
See Appendix.
\end{proof}
\end{proposition}

We consider the distribution of couples' wealth. Under random matching
(see the Appendix) this has mean \(0\) and variance \(2S^2\). Under SGAM, the variance is:
\[
V(x_{2}+y_{2}) = 2S^{2} + 2C^{2} \ge 2S^{2} 
\]
with strict inequality if \(a < 1\) or \(\sigma > 0\). The variance is decreasing in
\(a\) and equals \(4S^2\) if \(a = 0\). Thus, SGAM increases cross-sectional inequality,
but less so than pure matching on wealth.

\subsection{Children}\label{children}

All couples have the same number of children. A child's characteristics are
given by:
\begin{eqnarray}
x_{1}^{\prime } &=&\frac{\tau }{2}\left( x_{1}+y_{1}\right) +\varepsilon
\label{Chil} \\
x_{2}^{\prime } &=&\frac{\theta }{2}\left( x_{2}+y_{2}\right) +\eta 
\nonumber
\end{eqnarray}
where \(x\) and \(y\) are the child's parents, and \(\varepsilon\) and \(\eta\) are
independent normal random shocks with mean \(0\) and variance \(1\).

Parameter \(\tau \approx 1\) reflects genetic inheritance. Under standard biological
assumptions \(\tau = 1\) and characteristics show no regression to the mean. In our
model this leads the variance of \(x_1\) to grow without limit over generations.
In reality, we expect \(\tau < 1\) because very extreme characteristics are
selected against, a process known as stabilizing selection
(Schmalhausen 1949; Sanjak et al. 2018).

Parameter \(\theta \in [0, 1]\) reflects inheritance of SES. Unlike \(\tau\) it may vary
between societies. \(\theta\) is high when there is high intergenerational
transmission of SES. Among other things, \(\theta\) captures social and economic
institutions
that affect this intergenerational transmission, from taxation and
public education to hereditary nobility. If we interpret \(x_2\) narrowly as
wealth, \(1 - \theta\) can be thought of as the rate of inheritance tax.

For the time being, we assume that a person's genetic endowment has no impact on
their SES. Technically, thus, \(x_{2}^\prime\) does not directly depend on
\(x_{1}^\prime\). In a meritocratic society we would expect adult SES to partly depend
on genetics. We show that even absent meritocracy, correlations between \(x_1^\prime\) and
\(x_{2}^\prime\) can arise. In the Appendix, we relax this assumption and allow
meritocracy.

We can now calculate the covariance matrix for
\(x^\prime = (x_{1}^{\prime }, x_{2}^{\prime })\) under SGAM as:

\begin{eqnarray}
\mathbb{C} &=&\left( 
\begin{array}{cc}
1 & 0 \\ 
0 & 1%
\end{array}%
\right) +\left( 
\begin{array}{cccc}
\frac{\tau }{2} & 0 & \frac{\tau }{2} & 0 \\ 
0 & \frac{\theta }{2} & 0 & \frac{\theta }{2}%
\end{array}%
\right) \allowbreak \left( 
\begin{array}{cccc}
s^{2} & \sigma  & A^{2} & AC \\ 
\sigma  & S^{2} & AC & C^{2} \\ 
A^{2} & AC & s^{2} & \sigma  \\ 
AC & C^{2} & \sigma  & S^{2}%
\end{array}%
\right) \allowbreak \left( 
\begin{array}{cc}
\frac{1}{2}\tau  & 0 \\ 
0 & \frac{1}{2}\theta  \\ 
\frac{1}{2}\tau  & 0 \\ 
0 & \frac{1}{2}\theta 
\end{array}%
\right) \nonumber \\
&=&\left( 
\begin{array}{cc}
\frac{1}{2}A^{2}\tau ^{2}+\frac{1}{2}s^{2}\tau ^{2}+1 & \frac{1}{2}\theta
\sigma \tau +\frac{1}{2}AC\theta \tau  \\ 
\frac{1}{2}\theta \sigma \tau +\frac{1}{2}AC\theta \tau  & \frac{1}{2}%
C^{2}\theta ^{2}+\frac{1}{2}S^{2}\theta ^{2}+1%
\end{array}%
\right) \allowbreak \label{cov-children-SGAM}
\end{eqnarray}
\newline

We now explore two issues. First, under SGAM, genetic characteristics are no
longer exogenous; because of assortative matching, they are (partly) socially
determined. In particular, even if genetics and SES are uncorrelated among
parents, the expected genetic endowment of the child is positively related to
parental SES. Second, as a result, in the long run a correlation appears between
traits; that is, high SES people inherit genes that are attractive in marriage
markets.

Regarding point 1, we compute the expected genetic characteristic of the
child, conditional on parental wealth:
\[
\mathbb{E}\left[ \frac{\tau }{2}\left( x_{1}+y_{1}\right) +\varepsilon \mid
x_{2}=v,y_{2}=w\right] 
\]
Given the symmetry of the model, this conditional
expectation only depends on the parents' total wealth, i.e.~\(v+w\).

\begin{claim}\label{claim-E-children-SGAM}
Under SGAM, if $\sigma = 0$ (i.e. genetics and
SES are uncorrelated for the parents), then the expected genetic endowment of
the children is linearly increasing in parental SES. The relationship
increases with the ratio of genetic variance to SES variance, is zero for 
$a = 0$ or $a = 1$, and is highest for intermediate values of $a$.
\end{claim}

Next, we study the correlation between children's traits 1 and 2 as a function
of \(\sigma\), the covariance of parents' traits.

\begin{claim}\label{claim-corr-children-SGAM}
Under SGAM, if genetics and SES are uncorrelated for the parents, then they
are positively correlated for the children so long as $0 < a < 1$. The 
correlation is increasing in $\theta$.
\end{claim}

Whether characteristics are more or less correlated for children than for
parents depends on whether the initial correlation between parents'
characteristics is larger or smaller than the asymptotic one, derived below.

Figure \ref{fig:pic-intuition} shows the intuition behind the model. Parents
match on downward-sloping attractiveness isoquants given by \(a x_1 + (1-a) x_2 = u\).
Their children are in between them on both dimensions. This compresses the
distribution along the attractiveness isoquants, which leads to a
positive correlation between genetics and SES. The correlation between \(x^\prime_1\)
and \(x^\prime_2\) is 0 when \(a = 0\) or \(a = 1\), because then spouses don't trade off
SES for genes. It is highest for intermediate values of \(a\).

\begin{figure}

{\centering \subfloat[Parents\label{fig:pic-intuition-1}]{\includegraphics{model-section_files/figure-latex/pic-intuition-1} }\subfloat[Children\label{fig:pic-intuition-2}]{\includegraphics{model-section_files/figure-latex/pic-intuition-2} }

}

\caption{Theory. The shaded area is the population distribution. Parents (solid circles) match along attractiveness isoquants (dotted lines). Children (hollow circles) are between them. As a result, the children's distribution is squeezed along attractiveness isoquants, and $x_1$ and $x_2$ become associated.}\label{fig:pic-intuition}
\end{figure}

These results show that SGAM can lead to a genes-SES gradient, i.e.~a positive
correlation between genes and SES. Also, the strength of the genes-SES
correlation is affected by economic institutions, insofar as these are captured
in \(\theta\). When \(\theta\) is high, the genes-SES correlation is high too.

We now consider the asymptotic distribution of \(x_1\) and \(x_2\) when the matching
process is repeated over many generations. As we would expect, our main results
continue to hold.

\begin{proposition}\label{prop-asymptotics-SGAM}
Under SGAM, for $\theta <1$ and $\tau <1$, the dynamics converge to a stationary
distribution that is normal with mean zero and covariance matrix 
\[
\mathbb{C}\left(\begin{array}{cc}
x_1  \\
x_2
\end{array}
\right)
=
\left(
\begin{array}{cc}
\bar{s}^{2} & \bar{\sigma} \\ 
\bar{\sigma} & \bar{S}^{2}%
\end{array}%
\right) 
\]

where $\bar{s}, \bar{S}, \bar{\sigma}$ are functions of the parameters ($a, \theta, \tau$).
Moreover, the asymptotic correlation between characteristics, 
$corr = \bar{\sigma}/\bar{s}\bar{S}$, is non-negative, positive for $0 < a < 1$, 
increasing in $\theta$ and increasing then decreasing in $a$. 
The coefficient of parents' wealth on children's genetics is also positive 
for $0 < a < 1$.

For $\theta = 1$, the dynamics diverge and $\bar{S}^{2}$ goes to $+\infty$; for 
$\tau = 1$, the dynamics diverges and $\bar{s}^{2}$ goes to $+\infty$.

\end{proposition}

\begin{figure}

{\centering \includegraphics[width=6in]{model-section_files/figure-latex/pic-asymptotic-corr-1} 

}

\caption{Long-run squared correlation between genetics $x_1$ and SES $x_2$ (heritability of SES), by weight of genetics in spouse matching ($a$) and strength of direct transmission of SES ($\theta$). $\tau$ = 0.95.}\label{fig:pic-asymptotic-corr}
\end{figure}

Figure \ref{fig:pic-asymptotic-corr} plots the asymptotic squared correlation
between \(x_1\) and \(x_2\). In genetic terms, this is the heritability of SES. It is
maximized for intermediate levels of \(a\). Note that both \(\bar{S}^{2}\) and
\(\bar{\sigma}\), as well as the correlation between characteristics and the
conditional expectation of genetics given wealth, are increasing in \(\theta\).
Higher transmission of SES increases inequality of SES in equilibrium
(unsurprisingly), but also the heritability of SES.

\subsection{Extensions}\label{extensions}

In modern meritocracies, people's adult SES depends not just
on their parents' social status and on chance, but also on their own effort and
skills, which might be related to their genetics. So let

\begin{align}
x^\prime_1 &= \tau \frac{x_{1} + y_{1}}{2} + \varepsilon   \nonumber \\
x^\prime_2 &= \gamma x^\prime_1 + \theta \frac{x_{2}+y_{2}}{2}+\eta \label{eqn-gamma}
\end{align}

where \(\gamma > 0\) represents the effect of own genetics on own SES. \(\gamma\)
can be thought of as measuring meritocratic mobility, but alternatively as the
degree of overlap between genes that matter in marriage and labour
markets. The basic result continues to hold, and also, \(\gamma\)
increases the correlation between genes and SES.

\begin{proposition}\label{prop-gamma}
Under SGAM and equation \eqref{eqn-gamma}, if genetics and SES are 
uncorrelated for the parents, then they
are positively correlated for the children so long as $0 < a < 1$ or
$\gamma > 0$. The correlation is increasing in $\gamma$. Also, so 
long as $\gamma > 0$ and either $0 < a < 1$ or $\sigma > 0$, the coefficient 
of parents' wealth on children's wealth exceeds $\theta$.
\end{proposition}

With this extension, we can examine how intergenerational transmission
of inequality responds to policy shocks under different models of assortative
mating. Figure \ref{fig:pic-simulate-mao} simulates different generations of
families under a policy shock which reduces \(\theta\) and \(\gamma\) to 0 for
one generation. We group families into quartiles by their SES on the eve of the
shock, and track the quartiles' mean SES over generations. The left panel shows
pure genetic assortative mating (\(a = 1\)). Since couples only match on genetics,
although the shock temporarily reduces SES differences to zero, afterwards
families return to where they would have been without the shock
(dashed lines). The right panel shows SGAM with \(a = 0.3\). Here, the shock
reduces the mating value of previously high-SES individuals, and as a result,
their SES never recovers to its no-shock value.

\begin{figure}
\centering
\pandocbounded{\includegraphics[keepaspectratio,alt={\label{fig:pic-simulate-mao}A simulated policy shock under GAM (a = 1) and SGAM (a = 0.3). Lines show mean SES in each quartile of families, grouped by SES in the generation before the shock. In generation 0 , \textbackslash theta and \textbackslash gamma are reduced to zero. Thereafter they return to their previous values (\textbackslash theta = 0.8, \textbackslash gamma = 0.1. Under GAM, family SES returns to the no-shock counterfactual (dashed lines). Under SGAM it never does so.}]{model-section_files/model-section_files/figure-latex/pic-simulate-mao-1.pdf}}
\caption{\label{fig:pic-simulate-mao}A simulated policy shock under GAM (\(a\) = 1) and SGAM (\(a\) = 0.3). Lines show mean SES in each quartile of families, grouped by SES in the generation before the shock. In generation 0 , \(\theta\) and \(\gamma\) are reduced to zero. Thereafter they return to their previous values (\(\theta = 0.8, \gamma = 0.1\). Under GAM, family SES returns to the no-shock counterfactual (dashed lines). Under SGAM it never does so.}
\end{figure}

In the Appendix we consider further extensions. We allow men
and women to have different values of \(a\), i.e.~different
weights on genetics and SES in the attractiveness function. Lastly, we
show that our basic result is robust to allowing non-normal distributions of
\(x_1\) and \(x_2\) and non-linear attractiveness functions.

Other extensions are possible. We assumed that all couples have the same number
of children. If fertility increased with \(x_1\) or \(x_2\), we would expect this to
reduce the variance of traits in the children's generation and possibly also
their covariance. Here, matching preferences, summarized by the \(a\) parameter,
are exogenous; it would be natural to model \(a\) as an equilibrium outcome.
For example, if parents care about their children's wealth, \(a\) might decrease
in \(\theta\). Indeed, below we see suggestive evidence that
income is valued differently in Norwegian and British marriage markets. Finally,
if different individuals have varying matching preferences
(values of \(a\)), the model becomes hard to work with because there is no
agreed-upon ranking of men and women, but computations using the
Gale-Shapley matching algorithm shows that the basic insights of SGAM still hold.\footnote{See \url{https://github.com/hughjonesd/trading-status-genetics/blob/main/gale-shapley-multiple-a.R}.}

\subsection{Discussion}\label{discussion}

The meanings of both social status, and ``good genes'' in the
marriage market, are likely to vary across societies. Social status could
encompass variables like social class or caste; ethnic identity in
``ranked'' ethnic systems; or in modern societies, SES, including wealth, income
and occupation. Regarding genetics, standards of physical attractiveness, and
other genetically-influenced characteristics which make someone a ``good match'',
vary across societies and over time. The central prediction of the model is that
whatever those characteristics, in the long run they will become correlated with
SES.

Recent empirical work shows high persistence of SES over time, in particular at
the top. Clark (2023) argues that this could be explained by unobserved
genetic variation.
Proposition \ref{prop-gamma} shows that if genes affect own wealth directly,
under assortative mating, the regression coefficient of parents' wealth on own
wealth exceeds the ``direct'' coefficient \(\theta\), because parents' wealth
correlates with parents' genetics and via that with own wealth. Thus,
regressions of wealth on wealth may include the effect of unobserved genetic
variation. This may be a confound due to pre-existing gene-SES correlation (if
\(\sigma > 0\)). But under SGAM (\(0 < a < 1\)) it can also be a mediating
variable, since changes in someone's wealth may indeed affect the identity of
their spouse, hence the genetics of their offspring, and from that their
offspring's adult wealth.

The converse also holds: regressions of children's characteristics on their
genetics alone risk overestimating the effect of genetics, by confounding it
with the effects of correlated socio-economic status. Recent work in genetics
has shown this. Polygenic scores for educational attainment have smaller effects
in between-sibling regressions, where between-family variation in SES is
partialled out and where genetic variants are guaranteed to be randomly
allocated, than in regressions which pool the whole sample ({Howe et al.} 2022).
Parents' genetic variants which are \emph{not} passed on to children predict
children's characteristics, partly due to social stratification in the geneticists'
sense of non-random mating ({Kong et al.} 2018; Young 2023).

The model predicts variation in the effects of SGAM. In particular, in
``caste societies'' where there is complete endogamy within social status
groups, there is no scope for SGAM, because marriage partners do not trade off
genetics for social status (\(a = 0\)). Also, the association between genes and
SES is increased by the institutional variable \(\theta\), which captures
intergenerational persistence of SES. This implies that policy has long-run
effects on biosocial structure: reducing \(\theta\) not only increases
intergenerational mobility, but reduces the correlation of genes with SES, and
hence the bias of what Harden (2021) calls the ``genetic lottery''.
Conversely, reforms that increase meritocracy (\(\gamma\), Proposition
\ref{prop-gamma}) may strengthen the genes-SES gradient. This poses the problem
raised by Young (1958) and more recently Markovits (2019): meritocracy
may be self-limiting or even self-undermining.

\protect\phantomsection\label{refs}
\begin{CSLReferences}{1}{1}
\bibitem[\citeproctext]{ref-clark2023inheritance}
Clark, Gregory. 2023. {``The Inheritance of Social Status: England, 1600 to 2022.''} \emph{Proceedings of the National Academy of Sciences} 120 (27): e2300926120.

\bibitem[\citeproctext]{ref-harden2021genetic}
Harden, Kathryn Paige. 2021. \emph{The Genetic Lottery: Why DNA Matters for Social Equality}. Princeton University Press.

\bibitem[\citeproctext]{ref-howe2022within}
{Howe, Laurence J, Michel G Nivard, Tim T Morris, et al.} 2022. {``Within-Sibship Genome-Wide Association Analyses Decrease Bias in Estimates of Direct Genetic Effects.''} \emph{Nature Genetics} 54 (5): 581--92.

\bibitem[\citeproctext]{ref-kong2018nature}
{Kong, Augustine, Gudmar Thorleifsson, Michael L Frigge, et al.} 2018. {``The Nature of Nurture: Effects of Parental Genotypes.''} \emph{Science} 359 (6374): 424--28.

\bibitem[\citeproctext]{ref-markovits2019meritocracy}
Markovits, Daniel. 2019. \emph{The Meritocracy Trap}. Penguin UK.

\bibitem[\citeproctext]{ref-sanjak2018evidence}
Sanjak, Jaleal S, Julia Sidorenko, Matthew R Robinson, Kevin R Thornton, and Peter M Visscher. 2018. {``Evidence of Directional and Stabilizing Selection in Contemporary Humans.''} \emph{Proceedings of the National Academy of Sciences} 115 (1): 151--56.

\bibitem[\citeproctext]{ref-schmalhausen1949factors}
Schmalhausen, Ivan Ivanovich. 1949. \emph{Factors of Evolution: The Theory of Stabilizing Selection.}

\bibitem[\citeproctext]{ref-young2023estimation}
Young, Alexander Strudwick. 2023. {``Estimation of Indirect Genetic Effects and Heritability Under Assortative Mating.''} \emph{bioRxiv}, 2023--07.

\bibitem[\citeproctext]{ref-young1958rise}
Young, M. 1958. \emph{The Rise of the Meritocracy}. Thames; Hudson.

\end{CSLReferences}

\end{document}
